\section{Methods}

\subsection{Model description}

\subsubsection{Amphibian DEB model}
%
The DEB-TKTD model is based on DEBkiss \citep{jagerDEBkissQuestSimplest2013} 
and has been extended to account for amphibian metamorphosis \citep{hansulDynamicModellingAmphibian}. 
In short, the model includes a metamorphic reserve compartment $E^{mt}$, 
which is defined as the amount of biomass that is used during metamorphosis to fuel 
maintenance costs and maturation. 
The parameter $\gamma$ determines the allocation of assimilated resources into the metamorphic reserve 
compartment. 
Under \textit{ad libitum} feeding conditions and at equilibrium, 
$\gamma$ is equal to the fraction of biomass that is reserve.\\
During metamorphosis, feeding rates decline and are assumed to reach 0 with the completion of metamorphosis. 
The decline in the size-specific feeding rate is coupled to the depletion of metamorphic reserves:
%
\begin{equation}
    \{ \dot{I} \}_{max}^{mt} = \{ \dot{I} \}_{max}^{lrv} \dfrac{E^{mt}}{E^{mt}_{max}}  
    \label{eq:dI_mt}
\end{equation}
%
Where  $\{ \dot{I} \}_{max}^{lrv}$ is the maximum size-spefific ingestion rate for larvae,
$\{ \dot{I} \}_{max}^{mt}$ is the current maximum size-spefific ingestion rate for metamorphs 
and $E^{mt}_{max}$ is the internally tracked maximum reserve level, 
reached at metamorphic climax.\\
By definition, metamorphosis is completed when the metamorphic reserve is depleted. 
The full dynamics of the model as derived by \citet{hansulDynamicModellingAmphibian} are given by:
\begin{align}
    \dot{X}^{emb} &=
    \begin{cases}
        -\dot{I} & \text{if embryo} \\ 
        0 & \text{otherwise}
    \end{cases} \\
    \dot{X} &= 
    \begin{cases}
        0 & \text{if embryo} \\ 
        \dot{I} & \text{otherwise}
    \end{cases}\\
    \dot{I} &= 
    \begin{cases}
        0 & \text{if embryo} \\
        \{ \dot{I} \}_{\text{max}}\, f([X])\, S^{2/3} & \text{otherwise}
    \end{cases} \\
    \dot{A} &= y_A \eta_{IA} \dot{I} \\ 
    \dot{M} &= y_M k_M S \\ 
    \dot{J} &= y_M k_J H \\
    \dot{S} &= 
    \begin{cases}
        y_G \eta_{AS} (1-\gamma) (\kappa \dot{A} - k_M S) & \text{if larva} \\ 
        y_G \eta_{AS} \dot{A} & \text{if metamorph} \\
        y_G \eta_{AS} (\kappa \dot{A} - \dot{M}) & \text{otherwise}
    \end{cases} \\
    \dot{E}^{mt} &= 
    \begin{cases}
        y_G \eta_{AS} \gamma (\kappa \dot{A} - k_M S) & \text{if larva} \\
        -(\dot{M} + \dot{J} + \dot{H}) & \text{if metamorph} \\ 
        0 & \text{otherwise}
    \end{cases} \\
    \dot{H} &= 
    \begin{cases} 
        0 & \text{if adult} \\
        (1-\kappa) \dot{A} - \dot{J} & \text{otherwise} \\
    \end{cases} \\
    \dot{R} &= 
    \begin{cases} 
        y_R \eta_{AR} \left[ (1-\kappa) \dot{A} - \dot{J}\right] & \text{if adult}\\
        0 & \text{otherwise}
    \end{cases}
    \label{eq:fullmodel}
\end{align}
%
with the following rules for life stage transitions: 
%
\begin{align}
    embryo &= 
    \begin{cases}
        \text{true} & \text{if } X_{emb}^{int} > 0 \\
        0 & \text{otherwise}
    \end{cases} \\
    juvenile &= 
    \begin{cases}
        \text{true} & \text{if } X_{emb}^{int} = 0 \land H < H^p \\ 
        0 & \text{otherwise}
    \end{cases} \\
    adult &= 
    \begin{cases}
        \text{true} & \text{if } H \geq H^p \\ 
        0 & \text{otherwise}
    \end{cases}
    \label{eq:lifestages}
\end{align}%\todo{update life stage definitions}
%
The parameters and state variables are listed in Table \ref{tab:all_params}.
The values $y_G$, $y_M$, $y_A$ and $y_{\kappa}$ 
are intermediate quantitites which account for the effects of stressors.
{\begin{table}[ht]
    \caption{
        Model parameters with default values. 
        Note that parameters can in addition be life stage-specific, indicated by optional superscripts.
    }
    \label{tab:all_params}
    \begin{tabularx}{\textwidth}{lll}
        \hline
        Quantity & Unit & Description  \\ 
        \hline
        \multicolumn{3}{l}{\textbf{DEB Parameters}}\\
        \hline
        $Z$ & $-$ & Mass-based zoom factor \\ 
        $X^{emb}_{int}$ & $mg$ & Initial mass of vitellus, approximated as dry mass of an egg \\
        $K_X$ & $mg\ L^{-1}$  & Half saturation for food uptake \\
        $\{\dot{I}_{max}\}$  & $mg\ mg^{-2/3}\ d^{-1}$ \\
        $\eta_{IA}$ & $-$ & Assimilation efficiency \\
        $\eta_{AS}$ & $-$ & Growth efficiency \\
        $\eta_{AR}$ & $-$ & Reproduction efficiency \\
        $\kappa$ & $-$ & Allocation fraction to soma \\
        $\gamma$ & $-$ & Allocation to metamorphic reserves \\    
        $k_M$ & $d^{-1}$ & Somatic maintenance rate constant \\
        $k_J$ & $d^{-1}$ & Maturity maintenance rate constant \\
        $H^j$ & $mg$ & Maturity at metamorphosis (approximately GS 42) \\
        $H^p$ & $mg$ & Maturity at puberty \\
        \hline
        \multicolumn{3}{l}{\textbf{DEB state variables \& intermediate states}}\\
        \hline
        \multicolumn{3}{l}{\textbf{TKTD parameters}}\\
        \hline
        \multicolumn{3}{l}{\textbf{TKTD state variables \& intermediate states}}\\
    \end{tabularx}
\end{table}}%\todo{add state variables}
%
Additionally, the zoom factor $Z$, calculated as the ratio between the maximum structural masses of two organisms, 
is used to induce variability between individuals as described by \citet{hansulDynamicModellingAmphibian}. 
This leads to variability in growth rates, maximum size, time to reach metamorphosis, 
time to complete metamorphosis , 
and body mass at the beginning and end of metamorphosis, respectively.\\
DEB parameters, including $\sigma_{Z}$, have previously been inferred for \textit{Discoglossus galganoi}. 
Parameter values used in this study are taken from \citet{hansulDynamicModellingAmphibian} and reported in Appendix. 
%
\subsubsection{Toxicokinetic-toxicodynamic model}
%
To model the effects of chemicals, we first applied a minimal component for the damage dynamics: 
%
\begin{equation}
    \dot{D}_{z,j} = k_D \left( C_{W,z}  - D_{z,j} \right)
    \label{eq:TK}
\end{equation}
%
Equation \ref{eq:TK} links the aqueous concentration $C_{W,z}$ of stressor $z$ to 
the scaled damage $D_{z,j}$, which is specific to the physiological mode of action (PMoA) $j$.
This formulation ignores possible feedbacks beween scaled damage and body size \citep{jagerRevisitingSimplifiedDEBtox2020}, 
in an attempt to deploy the simplest possible model. \\ 
For each possible PMoA $j$, the corresponding damage $D_{z,j}$ is translated to a relative response $y_{z,j}$. 
We assume a log-logistic relationship between damage and response, 
corresponding to the standard assumption used in ecotoxicological analysis and resulting in a continuously differentiable 
TKTD model. 
For PMoAs with a monotonically decreasing relationship, the response can be directly calculated from the survival function 
of the log-logistic distribution:
%
\begin{equation}
    y_{z,j} = \dfrac{1}{ 1 + \left( \dfrac{D_{z,j}} {e_{z,j}} \right) ^ {\beta_{z,j}} }
    \label{eq:TD_decrease}
\end{equation}
%
This function has two parameters, the sensitivity $e_{z,j}$ ($mg/L$) an the slope $\beta_{z,j}$ (dimensionless). 
In the present study, PMoAs with decreasing response are: 
Decrease in growth efficiency ($G$), 
decrease in assimilation efficiency ($A$) 
and decreased investment in growth/increased investment in maturation ($\kappa$).
For these PMoAs, $e$ is interpretable as a median effective damage, i.e. the damage level  
for which the corresponding flux is decreased by 50\%. \\
For PMoAs with a monotonically increasing relationship, we use an affine transformation of the cumulative hazard 
function of the log-logistic distribution:
%
\begin{equation}
    y_{z,j} = 1 - ln \left[ \frac{1}{ 1 + \left( \dfrac{D_{z,j}} {e_{z,j}} \right) ^ {\beta_{z,j}} } \right]
    \label{eq:TD_increase}
\end{equation}
%
In the present study, the only PMoA with an increasing relationship is an increase in maintenance costs ($M$).
The parameter $e_{z,j}$ now corresponds to the damage level at which a 1.7-fold increase in maintenance costs occurrs.  
Equations \ref{eq:TD_decrease} and \ref{eq:TD_increase} are defined so that the resulting relative response is a factor 
that can be inserted directly into the DEB equations. 
The intermediate variable \textit{stress} \citep{jagerRevisitingSimplifiedDEBtox2020}
can be skipped without consequences for the model dynamics. \\
%
For mixtures of stressors with different PMoAs, their effects can be applied independently. 
For stressors with different PMoAs, we assume that an independent action (IA) model applies, 
where the combined relative response $y_j$ is the product of the individual responses: 
%
\begin{equation}
    y_{j} = \prod\limits_{z=1}^{n} y_{z,j}
    \label{eq:IA}
\end{equation}
%
Equation \ref{eq:IA} yields the values $y_G$, $y_M$, $y_A$ and $y_{\kappa}$ used in equation \ref{eq:fullmodel}.
The alternative to IA is a damage addition (DA) model. In the context of DEB-TKTD models, DA requires the estimation 
of an additional weight parameter for each substance. This means that the data is either fitted to the mixture data - 
which would contradict the goals of the present study - or that all TKTD parameters have to be estimated simultaneously 
from the combined single-stressor data. The assumption of IA is therefore a plausible and practical starting point.
%
\subsubsection{Parameter inference}
%
TKTD parameters were inferred from single-stressor effect data of 2,4-D and Flupyradifurone (product Sivanto, Bayer, Germany), respectively 
(Gonzalez-Lopez et al., unpublished).\\ 
The calibration data consisted of time-resolved larval wet mass, 
the time to reach Gosner stage (GS) 42 and the wet mass at GS 42. 
The time to reach GS 42 was used as the time-resolved fraction 
of tadpoles, calculated as the number of survivors which have not yet reached GS 42, 
relative to the total number of surivors at the given time point.
An slower decrease in this value over time therefore indicates delayed metamorphosis, 
unrelated to survival.  
The highest tested 2,4-D treatment (100 $mg/L$) was omitted from the calibration, 
because it caused 100\% mortality quickly after the onset of exposure. 
The presented analysis exclusively deals with sublethal effects.\\
% 
For each stressor, we pre-selected plausible PMoAs and separately fitted a TKTD model based on each PMoA. 
The plausibility of PMoAs was evaluated by visual comparison.
%
All parameter inferences were done using population monte carlo approximate Bayesian computation 
(PMC-ABC \citep{beaumontAdaptiveApproximateBayesian2009}), 
using the Eculidean distance as distance function, which is a default choice in ABC \citep{prangleAdaptingABCDistance2017}. 
PMC-ABC is a practical choice because the model output was stochastic, 
and PMC-ABC is well-suited to deal with stochasticity in model outputs. 
Details on the PMC-ABC configuration are given in supporting information.\\
%
\subsubsection{Model cross-validation and PMoA selection}
%
To cross-validate the model, we attempted to predict the effects at GS 46 from effects on larvae up to GS 42. 
To do so, we used the point estimates of parameters inferred for PMoAs which plausibly reproduced the data during parameter inference. 
From the raw simulation output, i.e. the solution of the ODE system \ref{eq:fullmodel}, 
we extracted the observed metamorphosis traits, mimicking the fixed-stage design. \\ 
We also used this part of the data to narrow down the selection of the PMoA, 
by comparing the predictions made by different PMoAs visually. 
\subsection{Model validation with mixture data}
%
The model was validated by predicting the effect of the binary mixture of 2,4-D and Flupyradifurone from 
the single-stressor effects.
This involved two sets of predictions: 
\begin{enumerate}
    \item Predicting the effects of the binary mixture on larvae up to GS 42 from the corresponding single-stressor data.
    \item Predicting the effects of the binary mixture at GS 46 from the single-stressor data on larvae up to GS 42.
\end{enumerate}
%
The latter prediction combines the prediction of mixture effects with the extrapolation of effects across life stages.
%
To assess the model's capacity quantitatively, we caluclated the mean absolute percent error (MAPE) 
for the selected PMoA: 
\begin{equation}
    MAPE = 100 \cdot \left[ \sum\limits_{i=1}^{n} \left( \dfrac{\hat{y}_i - y_i}{y_i} \right) n^{-1}  \right]
    \label{eq:mape}
\end{equation}
%
$y_i$ is the $i$th observed value and $\hat{y}_i$ is the $i$th simulated value.
